\telubab{5}{BAB V}{CONCLUSION}{sec:5}

\subsection{Summary}\label{sec:5.1}

This thesis asked whether a search agent trained with reinforcement
learning can be rewarded for \emph{grounded} answers, rather than only
correct ones, without an LLM judge and on a student budget.

\begin{itemize}
\tightlist
\item
  \textbf{Setup.} A Qwen2.5-7B-Instruct agent with LoRA was trained with
  a GRPO-style group-normalised policy gradient on HotpotQA. It searched
  a closed pool of each question's 10 paragraphs with BM25, and answered
  with sentence-level citations that a rule could check against
  HotpotQA's gold supporting facts.
\item
  \textbf{Runs.} Sixteen runs were trained on one GPU for about \$49 in
  total: five reward and optimizer configurations on three seeds each,
  plus one diagnostic run.
\item
  \textbf{Evaluation.} Each run was evaluated on a frozen 300-question
  holdout, with seed-aware bootstrap confidence intervals.
\end{itemize}

\textbf{RQ1: Does GRPO training improve a 7B search agent over the
un-tuned model?} Yes, mostly through protocol learning. Every method
gains about +20 EM points over the un-tuned model (22.3\%), and every
seed-aware confidence interval excludes zero. Most of the gain comes
from producing valid actions. The un-tuned model fails 130 of 300
episodes, while trained runs fail at most 22. EM on completed episodes
rises only from 39.4\% to 45.2\% (R1).

\textbf{RQ2: Does VEGR prevent the shortcuts of answer-only reward
without losing accuracy?} Yes, within this setting.

\begin{itemize}
\tightlist
\item
  \textbf{Accuracy.} The answer-only reward and the full VEGR bundle
  reach the same accuracy (43.6\% vs.~42.6\% EM; +1.0, seed-aware CI
  {[}−3.6, +5.4{]}).
\item
  \textbf{Grounding.} VEGR's precision-grounded EM is \textbf{24.4
  points higher} (CI {[}+8.8, +37.3{]}).
\item
  \textbf{Citation loss under the answer-only reward.} Two of three
  answer-only seeds stopped citing. The training rollouts explain why:
  under the answer-only reward, an invalid citation costs −1 and a
  correct one earns nothing, and early in training about half of all
  citation attempts were invalid. The third seed kept citing but drifted
  into repeated searches, which that reward neither punishes nor
  rewards.
\item
  \textbf{What the comparison changes.} R1 differs from R0 in six reward
  terms and in its number of updates, so the result concerns the VEGR
  bundle as a whole.
\item
  \textbf{Correct, not complete.} VEGR makes citations correct but not
  complete. Its evidence coverage is no higher than the un-tuned
  model's, and its Joint-EM is low and varies between seeds.
\end{itemize}

\textbf{RQ3: How do reward variants and optimizer changes affect
grounding and search behaviour?} These results are exploratory, and none
of the variants is shown to improve on plain VEGR.

\begin{itemize}
\tightlist
\item
  Gated penalties (R2) and advantages without std normalisation (R4)
  gave inconclusive results. R2's apparent single-seed gain of +6 EM
  points did not replicate.
\item
  Doubling the coverage weight (R3) raised evidence coverage from 0.44
  to 0.57 and gave the best HotpotQA Joint-EM (+6.3 over R1, CI {[}+0.2,
  +13.1{]}). It also lowered precision-grounded EM by 8.0 points (CI
  {[}−16.1, −0.2{]}). Which method is better depends on the metric.
\item
  On the one seed tested, dynamic sampling (R0-dyn) raised answer-only
  accuracy by 8.0 EM points, together with twice the training questions,
  but did not prevent the loss of citations.
\end{itemize}

\subsection{Contributions}\label{sec:5.2}

\begin{enumerate}
\def\labelenumi{\arabic{enumi}.}
\tightlist
\item
  \textbf{VEGR}, a judge-free, sentence-level evidence reward that also
  enforces the action protocol. Its individual terms are not new; the
  contribution is the combination and its controlled evaluation.
\item
  \textbf{An explanation of how answer-only reward loses citations}, in
  a small, reproducible setting where the behaviour can be detected by
  rule. Invalid citations are penalised and correct ones are never
  rewarded. The thesis also shows that restoring the updates lost to
  zero-variance groups did not prevent the collapse on the seed tested.
\item
  \textbf{A controlled multi-seed comparison} with seed-aware confidence
  intervals. It shows how a single-seed result can mislead, and that
  grounding conclusions depend on the grounding metric used.
\item
  \textbf{A documented single-GPU pipeline.} One run costs about \$2.25.
  Splits, configurations, training logs, per-episode results and the
  analysis scripts are in the repository. The LoRA adapters of 15 of the
  16 runs are kept; the R1 seed-7 adapter was lost.
\end{enumerate}

\subsection{Limitations}\label{sec:5.3}

\begin{itemize}
\tightlist
\item
  \textbf{Setting.} Retrieval is BM25 over each question's own 10
  paragraphs, on one dataset, so the search problem is far easier than
  open-web search. The un-tuned model already retrieves 83\% of the gold
  documents.
\item
  \textbf{External validity.} VEGR needs sentence-level gold evidence,
  which few datasets provide, and it was tested on one dataset with one
  model size. The results say nothing yet about settings without such
  labels.
\item
  \textbf{Scale.} Training uses 256 questions per run for one epoch, on
  a 7B model with LoRA, and stops before rewards converge.
\item
  \textbf{Statistics.} Three seeds and 300 holdout questions resolve
  large effects, such as the grounding gap between R1 and R0. They
  cannot resolve EM differences of a few points, and the frequency of
  R0's collapse is essentially unknown (2 of 3 seeds, Clopper--Pearson
  95\% CI {[}0.09, 0.99{]}).
\item
  \textbf{Design.} The R1--R0 comparison changes the whole reward bundle
  and the number of updates. R2--R4 and R0-dyn are exploratory and
  uncorrected for multiple comparisons. R0-dyn has one seed and reached
  250 of its 256 target updates.
\item
  \textbf{Metric--reward overlap.} VEGR is trained on the same gold
  facts that the grounding metrics use. Precision-grounded EM is lenient
  about missing evidence, and no gold-independent attribution check was
  run.
\item
  \textbf{Algorithm.} The training update is a single on-policy step per
  question group, so PPO clipping never activated. The algorithm is a
  GRPO-style policy gradient, not full GRPO.
\item
  \textbf{Environment.} The environment accepts uncited answers, the
  citation separator collides with `\&' in some titles, and the reward
  weights were not tuned.
\end{itemize}

\subsection{Ethics and Broader Impact}\label{sec:5.4}

Answers with citations look trustworthy, and that appearance can lead
users to over-trust them. A citation that matches a gold sentence shows
that the agent pointed to relevant evidence. It does not show that the
answer was actually derived from that evidence, so precision-grounded EM
is not a measure of faithfulness. Even VEGR's best agents leave about a
quarter of their correct answers with at least one non-gold citation,
and they rarely cite the complete evidence. Systems built this way
should therefore present citations as pointers for checking, not as
proof.

The work uses only public data and models. HotpotQA is distributed under
the CC BY-SA 4.0 license. Qwen2.5-7B-Instruct is released by the Qwen
team under its own model license, which governs any use or
redistribution of the trained adapters. The total compute footprint was
modest: about \$49 of rented time on a single L40S GPU
(Table~\ref{tbl:3.2}). The findings are limited to one dataset with gold
evidence labels, one model and a closed retrieval pool, and should not
be read as evidence about open-web search agents.

\subsection{Future Work}\label{sec:5.5}

\begin{itemize}
\tightlist
\item
  \textbf{Separate the parts of the VEGR bundle.} Train VEGR without its
  evidence terms (coverage and precision), keeping the answer and search
  terms. This is the missing ablation needed to tell whether the
  grounding gain comes from rewarding evidence or from a denser reward
  with more updates. An evidence-terms-only variant would complete the
  picture.
\item
  \textbf{Close the no-citation loophole and fix the citation format.}
  Require at least one citation, and use a separator that cannot occur
  in titles (Section~\ref{sec:4.7}). Then test whether answer-only
  training still loses citations, or finds a different shortcut.
\item
  \textbf{Reward complete evidence without losing precision.} Examples
  are a reward based on the F1 of the citation set, or a coverage term
  counted only for gold citations, instead of a larger weight on
  coverage alone (R3).
\item
  \textbf{Check attribution without the gold labels.} Use an NLI-based
  citation check (Gao et al., 2023) to test whether VEGR's gains hold
  under a measure it was not trained on.
\item
  \textbf{Add a prompt-only baseline.} Since most of RQ1's gain is
  protocol learning, compare with a few-shot prompt that shows worked
  examples of the action format.
\item
  \textbf{Step-level credit.} Reward individual searches by the
  information they add, as in StepSearch (Wang et al., 2025) and STAMP
  (Xu et al., 2026), rather than scoring only the final episode.
\item
  \textbf{Harder retrieval and more data.} Extend the setup to an open
  corpus, to other multi-hop datasets with gold evidence such as
  2WikiMultiHopQA and MuSiQue, and to larger training sets.
\item
  \textbf{More seeds for the diagnostic runs.} Run R0-dyn on all three
  seeds with the same number of training questions as R0, so that
  dynamic sampling can be separated from extra data.
\end{itemize}
